\section{INTRODUÇÃO}\label{sec:introducao}

Este estudo analisa como o federalismo fiscal brasileiro, ao repassar aos municípios recursos da União, por meio do FPM, e dos estados, por meio da cota-parte do ICMS, pode distorcer a leitura de quanto cada um depende, de fato, de recursos de fora. Nesse escopo, foram analisados dois municípios da Bahia: Barreiras, no oeste baiano, e Camaçari, na Região Metropolitana de Salvador. Os dois recebem exatamente o mesmo valor do FPM desde 2018, R\$~214,1 milhões cada um em 2025, mas têm economias opostas: Camaçari é sede de um polo industrial, e Barreiras vive dos serviços regionais e da agropecuária. Em Barreiras, as transferências financiam cerca de 80\% da receita corrente; em Camaçari, entre 56\% e 65\%. O estudo compara ainda esses casos com a Hungria, que, diferentemente do Brasil, é um Estado unitário e não garante aos governos locais fatias automáticas da arrecadação nacional.

Esse contraste interessa às finanças públicas porque toca dois problemas clássicos do federalismo fiscal. O primeiro é o efeito \textit{flypaper}. Pelo modelo do eleitor mediano, uma transferência incondicional deveria elevar o gasto local tanto quanto um aumento equivalente na renda dos moradores; na prática, eleva muito mais \citep{HinesThaler1995}. O segundo é a ilusão fiscal: quando parte do custo dos serviços é paga por terceiros ou por tributos pouco visíveis, os contribuintes subestimam esse custo e passam a demandar mais gasto \citep{DolleryWorthington1996}. A tese que serve de referência a este trabalho trata o primeiro fenômeno como uma forma do segundo. Ela encontra o efeito em estados e municípios brasileiros, mais forte onde é mais custoso arrecadar com tributos locais, o chamado custo marginal dos fundos públicos \citep{Ferreira2021}.

As duas ideias dependem de uma separação simples: o que é receita gerada pela economia local e o que é dinheiro vindo de fora. A contabilidade pública brasileira mistura as duas coisas. A cota-parte do ICMS é registrada como transferência do estado, mas pelo menos 65\% dela é repartida conforme o valor adicionado em cada município \citep{CF1988art158}. Num município industrial, essa ``transferência'' sobe e desce com a atividade das próprias empresas; num município sem base produtiva relevante, funciona como um repasse redistributivo. Mesmo assim, a contabilidade dá às duas situações o mesmo nome: dependência de transferências. A distinção não é nova: a OCDE já separa quem controla a receita de quem a gera \citep{BlochligerKing2006}, e o Índice de Autonomia Local mede à parte o poder de tributar e conta a parcela incondicional dos repasses como autonomia de gasto \citep{Ladner2016}. O que este artigo examina são suas consequências em dois municípios que recebem exatamente o mesmo FPM.

O \textit{flypaper} e a ilusão fiscal ajudam a entender por que as transferências podem desestimular a arrecadação própria e esconder do eleitor o custo dos serviços. Há, porém, um passo anterior: saber o que é, de fato, transferência. Parte da renda gerada no município é tributada pelo estado e volta a ele como cota do ICMS, na proporção do que ali se produziu; outra parte da receita, como o FPM, vem da arrecadação nacional e é repartida pela população. A contabilidade trata as duas como dependência. O artigo pergunta, então, \textit{em que medida as transferências incondicionais se associam ao gasto de Barreiras e Camaçari mais do que a renda local, e como a dependência desses municípios se compara à dos governos locais da Hungria}. Sustenta-se que a dependência, como a contabilidade a registra, não acompanha a autonomia fiscal. No Brasil, ela ordena corretamente os municípios, mas superestima a dependência de recursos externos de quem tem a receita vinculada à própria economia; na Hungria, caiu enquanto os governos locais perdiam competências. Na parte econométrica, o gasto acompanha as transferências mais do que o PIB, mas, com dois municípios, não é possível separar esse padrão do ciclo comum das receitas.

A Hungria, segundo caso do estudo, percorreu o caminho inverso. A partir de 2011, com a nova Constituição e a nova lei de governos locais, o país recentralizou suas finanças locais em duas ondas \citep{Dobos2025}. Na primeira, o Estado assumiu as escolas, ao fim de 2012 \citep{Semjen2023}, e as dívidas acumuladas pelos municípios \citep{Vasvari2020571}. Desde 2013, passou também a financiar tarefas específicas com recursos carimbados, em vez de repartir receitas \citep{OCDE_UCLG2016}. Na segunda onda, a partir de 2019, os recursos discricionários passaram a favorecer prefeituras alinhadas ao governo e a faltar às de oposição \citep{Vasvari2024467, Kovarek2023}. Ao mesmo tempo, a contribuição de solidariedade, criada em 2017, passou a recolher ao orçamento nacional uma parte crescente do imposto local dos municípios que mais arrecadam \citep{CoE2024}. O resultado mostra o limite do indicador: os governos locais encolheram de 9,2\% para 5,4\% do PIB entre 2012 e 2025 e perderam poder, mas o peso das transferências em sua receita caiu. Pelo indicador usual, pareceriam mais autônomos. A comparação serve, portanto, para ver se a medida de dependência diz a verdade nos dois sentidos, e não para confrontar números de dois países: a lição húngara está menos no modelo unitário do que no indicador.

Para responder à pergunta, o artigo atualiza e corrige a base de dados do trabalho anterior sobre os mesmos casos \citep{DutraP1}, estendendo-a até 2025; compara a composição das receitas dos dois municípios; e estima a relação entre gasto, transferências e PIB, com testes de robustez. A principal contribuição é mostrar, com uma base corrigida e reprodutível, que, sob o mesmo FPM, o indicador convencional mais que dobra a dependência de recursos externos do município industrial, porque conta como transferência a receita gerada por sua própria indústria. A análise descritiva cobre de 2016 a 2025; as regressões vão de 2014 a 2023, último ano com PIB municipal. A Seção~\ref{sec:referencial} discute a literatura e apresenta o panorama bibliométrico; a Seção~\ref{sec:metodologia} descreve dados e métodos; a Seção~\ref{sec:resultados} apresenta os resultados; e a Seção~\ref{sec:conclusao} conclui.
